\usepackage[T1]{fontenc}
\usepackage{lmodern}
\usepackage{amsmath}

% we use biblatex, which is a more modern alternative to the traditional bibtex
% comment this to use bibtex instead
\usepackage[
    backend=biber,
    style=alphabetic,
    sorting=nyt,
    maxbibnames=99
]{biblatex}

% graphicx is needed to include images
\usepackage{graphicx}

% float is needed for [H] strict placement
\usepackage{float}

% xspace is used to add a space after a macro, if needed
\usepackage{xspace}

% acronym (https://ctan.org/pkg/acronym) is a package to expand acronyms after their first usage
\usepackage{acronym}

\usepackage{algorithm}
\usepackage{algpseudocode}

% the hyperref package provides links; with this configuration they are colored
% according to the scheme described in https://tex.stackexchange.com/a/525297
\usepackage{xcolor}
\usepackage[colorlinks]{hyperref}
\def\tmp#1#2#3{%
  \definecolor{Hy#1color}{#2}{#3}%
  \hypersetup{#1color=Hy#1color}}
\tmp{link}{HTML}{800006}
\tmp{cite}{HTML}{2E7E2A}
\tmp{file}{HTML}{131877}
\tmp{url} {HTML}{8A0087}
\tmp{menu}{HTML}{727500}
\tmp{run} {HTML}{137776}
\def\tmp#1#2{%
  \colorlet{Hy#1bordercolor}{Hy#1color#2}%
  \hypersetup{#1bordercolor=Hy#1bordercolor}}
\tmp{link}{!60!white}
\tmp{cite}{!60!white}
\tmp{file}{!60!white}
\tmp{url} {!60!white}
\tmp{menu}{!60!white}
\tmp{run} {!60!white}

% additional colors for diagrams and listings
\definecolor{amber}{RGB}{255, 191, 0}
% the booktabs package is used to create better tables; check an example of usage in results.tex
\usepackage{booktabs}
\usepackage{tabularx}
\usepackage{ragged2e}

%breaks long URLs better
\usepackage{xurl}

% macros
\newcommand*{\latex}{\LaTeX \xspace}
% Enable hyphenation and line breaking for typewriter font (\texttt / \ttfamily)
\makeatletter
\g@addto@macro\ttfamily{\hyphenchar\font=45\relax}
\makeatother
\renewcommand{\_}{\textunderscore\allowbreak}
\setlength{\emergencystretch}{2em}

\usepackage{hyperref}
\newcommand{\myhref}[2]{#2\footnote{\url{#1}}}

% graphics and diagrams
\usepackage{tikz}
\usetikzlibrary{arrows.meta, shapes.geometric, positioning}

% tracking changes using the todonotes package
\setlength {\marginparwidth }{2cm}
\usepackage{todonotes}
\newcommand{\note}[3]{\todo[inline,color=#2]{#1: #3}}
\newcommand{\md}[1]{\note{MD}{yellow}{#1}}

% result is something like "Related Work (Chapter 2)"
\newcommand{\fullref}[1]{\nameref{#1} (\Cref{#1})}

% used for the inparaenum environment
\usepackage{paralist}

% for code
\usepackage{listings}
\definecolor{listinggray}{gray}{0.9}
\definecolor{lbcolor}{rgb}{0.9,0.9,0.9}
\lstloadlanguages{C++,C,Java,Python,Bash,[x86masm]Assembler}

\lstdefinelanguage{CSharp}[Sharp]{C}{morekeywords={yield}}

\lstdefinelanguage{x86}[x86masm]{Assembler}{morekeywords={syscall,sysenter,section,global,bits,rel}}

\lstdefinelanguage{arm}{
  basicstyle=\ttfamily\small,
  alsoletter=-,
  keywords={MOV, ADD, SUB, SUBS, BNE, B, BL, LDR, SUBGT, SUBLT, CMP, MOVEQ, BLEQ, AND, ORR, EOR, BIC, MOVGT, CMPNE, MVN, CMN, TST, TEQ, ADC, SBC, RSB, RSC, LDMIA, STMIA, ADDLS, END, MOVS, SUBNES, MULNE, UMULL, UMLAL, SMULL, SMLAL, STMFD, LDMFD, PUSH, POP, STMDB, ANDS, BEQ, STR, LDRLS, SWI, SVC, BX, BGT, SUBLE, STRB, ASR, LSL, LSR, ROR, RRX, ORN, LDM, STM, ADDLE, SUBSNE, MOV32, MOVT},
  keywordstyle=\color{red}\bfseries,
  ndkeywords={.data, .text},
  ndkeywordstyle=\color{blue}\bfseries,
  identifierstyle=\color{black},
  sensitive=false,
  comment=[l]{@},
  commentstyle=\color[rgb]{0.133,0.545,0.133},
  stringstyle=\color[rgb]{0.627,0.126,0.941},
  morestring=[b]',
  morestring=[b]",
  frame=0
}

\lstdefinelanguage{JavaScript}{
  keywords={break, case, catch, continue, debugger, default, delete, do, else, false, finally, for, function, if, in, instanceof, new, null, return, switch, this, throw, true, try, typeof, var, void, while, with},
  morecomment=[l]{//},
  morecomment=[s]{/*}{*/},
  morestring=[b]',
  morestring=[b]",
  ndkeywords={class, export, boolean, throw, implements, import, this},
  keywordstyle=\color{blue}\bfseries,
  ndkeywordstyle=\color{darkgray}\bfseries,
  identifierstyle=\color{black},
  commentstyle=\color{purple}\ttfamily,
  stringstyle=\color{red}\ttfamily,
  sensitive=true
}

\lstdefinelanguage{yaml}{
  keywords={true,false,null,y,n,yes,no,True,False,TRUE,FALSE,NULL,on,off},
  keywordstyle=\color[rgb]{0,0,1}\bfseries,
  basicstyle=\ttfamily\footnotesize,
  sensitive=false,
  comment=[l]{\#},
  commentstyle=\color[rgb]{0.133,0.545,0.133},
  stringstyle=\color[rgb]{0.627,0.126,0.941},
  morestring=[b]',
  morestring=[b]",
  mathescape=false
}
\lstalias{YAML}{yaml}

\lstset{
  basicstyle=\ttfamily\small,
  language=C++,
  %backgroundcolor=\color{lbcolor},
  frame=single,
  commentstyle=\color[rgb]{0.133,0.545,0.133},
  keywordstyle=\color[rgb]{0,0,1},
  stringstyle=\color[rgb]{0.627,0.126,0.941},
  showtabs=false,
  showspaces=false,
  showstringspaces=false,
  columns=fullflexible,
  mathescape=true,
  keepspaces,
  breaklines=true,
  breakatwhitespace=false,
  upquote=true,
  tabsize=2,
  moredelim=[is][\color{red}]{`}{`},%
  moredelim=[is][\color{red}]{~red~}{~red~},%
  moredelim=[is][\color{green}]{~green~}{~green~},%
  moredelim=[is][\color{cyan}]{~cyan~}{~cyan~},%
  moredelim=[is][\color{gray}]{~gray~}{~gray~},%
  moredelim=[is][\color{blue}]{~blue~}{~blue~},%
  moredelim=[is][\color{blue}\bfseries]{/*blue*/}{/*blue*/},
  moredelim=[is][\color{red}\bfseries]{/*red*/}{/*red*/},
  moredelim=[is][\color{cyan}\bfseries]{/*cyan*/}{/*cyan*/},
  moredelim=[is][\bfseries]{/*bf*/}{/*bf*/},
  moredelim=[is][]{/***}{***/}
}


\usepackage{subcaption}

\usepackage[noabbrev,capitalise]{cleveref}
\AtBeginDocument{%
  \crefname{lstlisting}{code}{codes}%
  \Crefname{lstlisting}{Code}{Codes}%
  \crefname{listing}{code}{codes}%
  \Crefname{listing}{Code}{Codes}%
  \renewcommand{\lstlistingname}{Code}%
}

\renewcommand{\arraystretch}{1.2}
